\section*{Appendix}
\addcontentsline{toc}{section}{Appendix}
\label{sec:appendix}

\subsection*{A. Streaming Wire Protocol JSON Schema Specifications}

To ensure full interoperability across heterogeneous forensic microservices, our platform specifies strict JSON Schema definitions for all emitted streaming events.

\subsubsection*{1. Node Discovery Event Schema}
\begin{verbatim}
{
  "$schema": "http://json-schema.org/draft-07/schema#",
  "title": "NodeDiscoveredEvent",
  "type": "object",
  "properties": {
    "type": { "type": "string", "enum": ["node"] },
    "data": {
      "type": "object",
      "properties": {
        "id": { "type": "string", "pattern": "^(0x[a-f0-9]{40}|T[A-Za-z0-9]{33})$" },
        "label": { "type": "string" },
        "type": { "type": "string", "enum": ["eoa", "contract", "mixer", "bridge"] },
        "risk_score": { "type": "number", "minimum": 0, "maximum": 100 },
        "taint_ratio": { "type": "number", "minimum": 0, "maximum": 1 },
        "is_root": { "type": "boolean" },
        "tags": { "type": "array", "items": { "type": "string" } }
      },
      "required": ["id", "risk_score", "taint_ratio", "is_root"]
    }
  },
  "required": ["type", "data"]
}
\end{verbatim}

\subsubsection*{2. Wire Discovery Event Schema}
\begin{verbatim}
{
  "$schema": "http://json-schema.org/draft-07/schema#",
  "title": "WireDiscoveredEvent",
  "type": "object",
  "properties": {
    "type": { "type": "string", "enum": ["wire"] },
    "data": {
      "type": "object",
      "properties": {
        "source": { "type": "string" },
        "target": { "type": "string" },
        "amount": { "type": "number", "minimum": 0 },
        "token_symbol": { "type": "string" },
        "tx_hash": { "type": "string", "pattern": "^0x[a-f0-9]{64}$" },
        "priority_score": { "type": "number", "minimum": 0, "maximum": 1 },
        "timestamp": { "type": "integer" }
      },
      "required": ["source", "target", "amount", "token_symbol", "tx_hash"]
    }
  },
  "required": ["type", "data"]
}
\end{verbatim}

\subsection*{B. Supplementary Mathematical Proofs}

\begin{lemma}[Uniqueness of Canonical Address Representation]
Let $\mathcal{A}_{\text{raw}}$ be the set of arbitrary-case hexadecimal address strings. The normalization mapping:
\begin{equation}
    \eta(a) = \begin{cases}
    \text{lower}(a), & \text{if } a \in \{0\text{x}\} \times \{0\text{--}9, a\text{--}f, A\text{--}F\}^{40} \\
    a, & \text{otherwise}
    \end{cases}
\end{equation}
is a surjective projection onto the canonical address space $\mathcal{A}^*$, establishing an equivalence relation:
\begin{equation}
    a_1 \sim a_2 \iff \eta(a_1) = \eta(a_2)
\end{equation}
which preserves hash-table key consistency $\mathcal{O}(1)$ across heterogeneous runtime environments.
\end{lemma}

\begin{proof}
Hexadecimal equivalence is defined over value equality modulo $16$. The transformation $\text{lower}(c)$ maps uppercase glyphs $\{A\dots F\}$ onto $\{a\dots f\}$, which have identical numeric evaluation in base 16. Thus, for any two checksum variations $a_1$ and $a_2$ representing the same 160-bit integer, $\eta(a_1) = \eta(a_2)$. This ensures that dictionary lookups $H(\eta(a_1)) \equiv H(\eta(a_2))$, preventing duplicate node allocation and orphaned wires.
\end{proof}

\begin{lemma}[Non-Cyclic Traversal Invariance of Priority Queues]
Let $\mathcal{G}$ contain a directed cycle $\mathcal{C} = (u_1 \to u_2 \to \dots \to u_k \to u_1)$. Under Algorithm~\ref{alg:priority_search}, the traversal engine cannot enter an infinite loop.
\end{lemma}

\begin{proof}
Algorithm~\ref{alg:priority_search} maintains two invariant state sets: the visited node set $\mathcal{V}_{\text{visited}}$ and the admitted canvas edge set $\mathcal{E}_{\text{canvas}}$.
Upon the first traversal of edge $(u_k, u_1)$, node $u_1$ is inspected. Since $u_1 \in \mathcal{V}_{\text{visited}}$ (having been visited at the start of the cycle), line 27 of Algorithm~\ref{alg:priority_search} evaluates to false. Outgoing edges from $u_1$ are not re-fetched or re-enqueued. Furthermore, wire $(u_k, u_1)$ is added to $\mathcal{E}_{\text{canvas}}$ on line 15; any redundant edge pop skips expansion via line 13. Hence, the cycle is traversed at most once, and the queue strictly drains, guaranteeing termination.
\end{proof}

\subsection*{C. Framework Hyperparameter Reference Table}

Table~\ref{tab:hyperparameters} catalogs the complete hyperparameter configuration space utilized across empirical benchmarks and production deployments.

\begin{table}[h]
\centering
\caption{Framework Hyperparameter Configuration Reference.}
\label{tab:hyperparameters}
\begin{tabular}{llcl}
\toprule
\textbf{Parameter} & \textbf{Symbol} & \textbf{Default Value} & \textbf{Functional Description} \\
\midrule
Volume Weight & $\omega_1$ & 0.45 & Weight of logarithmic transacted amount in edge priority. \\
Taint Weight & $\omega_2$ & 0.25 & Weight of sender node scalar taint coefficient. \\
Velocity Weight & $\omega_3$ & 0.15 & Weight of temporal velocity decay factor. \\
Depth Weight & $\omega_4$ & 0.15 & Hop depth distance penalty weight. \\
Pruning Threshold & $\theta$ & 0.15 & Minimum priority score required for edge expansion. \\
Max Search Depth & $d_{\max}$ & 4 & Maximum graph hop distance from incident root. \\
Peel Volume Ratio & $\Omega_{\text{peel}}$ & 4.0 & Minimum asymmetry ratio between forward and peel wire. \\
Peel Dwell Limit & $\Delta T_{\text{peel\_max}}$ & $3,600\,\text{s}$ & Maximum dwell time between incoming and forward hop. \\
Smurf Min Degree & $K_{\text{smurf\_fan}}$ & 5 & Minimum out-degree to trigger smurfing evaluation. \\
Smurf Variance Bound & $\epsilon_{\text{variance}}$ & 0.35 & Maximum coefficient of variation across fan-out wires. \\
Stream Buffer Size & $M_{\text{queue}}$ & 500 & Maximum item capacity of async streaming event queue. \\
RPC Request Timeout & $T_{\text{rpc\_timeout}}$ & $2.0\,\text{s}$ & Network HTTP timeout before triggering benchmark fallback. \\
\bottomrule
\end{tabular}
\end{table}
